\subsection{Thiết kế PhoBERTIntentClassifier}

PhoBERTIntentClassifier là thành phần quan trọng trong thiết kế NLU tiếng Việt. Thành phần này nhận câu hỏi đã được chuẩn hóa và tách từ, sau đó sử dụng PhoBERT để tạo biểu diễn ngữ nghĩa toàn câu. Biểu diễn này được đưa qua lớp phân loại để dự đoán intent và confidence.

\begin{figure}[H]
    \centering
    \includegraphics[width=0.98\textwidth]{\detokenize{content/source/rasa/Cấu trúc module NLU và PhoBERTIntentClassifier.png}}
    \caption{Cấu trúc module NLU và PhoBERTIntentClassifier}
    \label{fig:rasa-nlu-phobert-classifier}
\end{figure}

Hình~\ref{fig:rasa-nlu-phobert-classifier} thể hiện đầy đủ các bước chính: preprocessor, VnCoreNLP, tokenizer, tensor input, PhoBERT encoder, pooling, classification head, softmax, ranking, confidence calibration và fallback gate. Đây là cấu trúc giúp Rasa hiểu câu hỏi ở mức ngữ nghĩa thay vì chỉ khớp từ khóa.

\begin{figure}[H]
    \centering
    \includegraphics[width=0.98\textwidth]{\detokenize{content/source/rasa/Kiến trúc xử lý của PhoBERT trong bài toán phân loại intent.png}}
    \caption{Kiến trúc xử lý của PhoBERT trong bài toán phân loại intent tiếng Việt}
    \label{fig:rasa-phobert-intent-architecture}
\end{figure}

Hình~\ref{fig:rasa-phobert-intent-architecture} đi sâu hơn vào cơ chế PhoBERT cho intent classification. Sau khi được token hóa, câu hỏi được đưa qua PhoBERT encoder để tạo vector ngữ cảnh. Vector này đi qua pooling và classification head, sau đó softmax tính xác suất cho từng intent. Intent có xác suất cao nhất được đề xuất cho Dialogue Management, miễn là vượt qua ngưỡng confidence.

Để làm rõ hơn vai trò của PhoBERTIntentClassifier, có thể xét một ví dụ cụ thể. Câu hỏi ban đầu của người dùng được chuẩn hóa, tách từ, mã hóa thành embedding, đưa qua tầng softmax để tính xác suất intent và cuối cùng trả về output cho Rasa gồm intent, confidence và ranking.

\begin{figure}[H]
    \centering
    \includegraphics[width=0.98\textwidth]{\detokenize{content/source/rasa/Ví dụ biến đổi câu hỏi qua PhoBERTIntentClassifier .png}}
    \caption{Ví dụ biến đổi câu hỏi qua PhoBERTIntentClassifier}
    \label{fig:rasa-phobert-question-transform}
\end{figure}

Hình~\ref{fig:rasa-phobert-question-transform} minh họa quá trình biến đổi một câu hỏi tiếng Việt qua PhoBERTIntentClassifier. Câu hỏi gốc trước hết được chuẩn hóa để giảm nhiễu ký tự và thống nhất chữ viết. Sau đó, văn bản được tách từ để tạo chuỗi token phù hợp với tiếng Việt. PhoBERT tiếp tục mã hóa chuỗi token thành biểu diễn ngữ nghĩa toàn câu. Tầng softmax tính xác suất cho từng intent, ví dụ intent hỏi điểm chuẩn đạt xác suất cao nhất. Kết quả cuối cùng được trả về cho Rasa dưới dạng intent, confidence và intent ranking để Dialogue Management lựa chọn action tiếp theo.

